% !TeX root = 微分几何学习笔记.tex
% ============================================================
%  封面：李思风（极简学术）—— 白底、居中标题、无装饰边框
%  参照《线性代数：原理与应用》（李思，清华大学）封面版式
%
%  可用字段（在主文件导言区 \newcommand 预定义即可覆盖）：
%    \notename      主标题（必填，主文件定义）
%    \noteauthor    作者（必填，主文件定义）
%    \notesubtitle  标题下方小字，默认 \today（日期/版本行）
%    \noteextra     底部单位/学院行，默认不显示
%
%  说明：中部线稿为可删装饰，整段 \begin{tikzpicture}...\end{tikzpicture}
%        删除后即为纯文字极简封面。
% ============================================================
\begin{titlepage}
	\thispagestyle{empty}
	\providecommand{\notesubtitle}{\today}
	\providecommand{\noteextra}{}
	\begin{center}
		\vspace*{4.7cm}

		% ---- 主标题 ----
		{\bfseries\songti\fontsize{24.8}{30}\selectfont \notename\par}

		\vspace{4.1cm}

		% ---- 日期 / 版本行 ----
		{\songti\fontsize{14.4}{18}\selectfont （\notesubtitle）\par}

		\vspace{3.4cm}

		% ---- 手绘风极简线稿（可整段删除）----
		\begin{tikzpicture}[line width=0.8pt, line cap=round]
			\draw[-{Stealth}, shorten >=1pt] (-2.1,0) -- (2.1,0);
			\draw[-{Stealth}, shorten >=1pt] (0,-0.9) -- (0,1.9);
			\draw[domain=-1.9:1.9, samples=80, smooth]
				plot (\x, {0.75*sin(deg(1.35*\x))+0.35});
			\draw[dashed, line width=0.6pt] (-1.25,-0.85) -- (-1.25,1.25);
			\node[font=\footnotesize] at (2.25,0.12) {$x$};
			\node[font=\footnotesize] at (0.14,2.0) {$y$};
			\node[font=\footnotesize] at (-1.45,1.32) {$f(x)$};
		\end{tikzpicture}

		\vspace{4.3cm}

		% ---- 作者（楷体）----
		{\kaishu\fontsize{24.8}{30}\selectfont \noteauthor\par}

		\vspace{1.8cm}

		% ---- 单位 / 学院（未定义则不显示）----
		{\songti\fontsize{20.7}{26}\selectfont \noteextra\par}

		\vspace*{1.6cm}
	\end{center}
\end{titlepage}
